\subsection{1. 引言}\label{introduction}

组合（composition）---用更简单的部件组装复杂系统---是软件工程的一项基础性原则 {[}1{]}。传统上，组合是静态的：函数调用、模块导入和类继承在编译期即被解析，并在整个执行过程中保持不变。然而，现代软件日益需要动态组合（dynamic composition），即组件在运行时被加载、卸载和重新配置。插件架构 {[}2{]} 与自演化智能体框架都要求系统能够在运行过程中安全地添加和移除功能；然而，当前的实践却转而依赖粗粒度机制 {[}3{]}，这类机制只能通过重启（并丢弃运行时状态）来完成重配置。尽管动态组合的实践重要性日益增长，但与静态组合已有的丰富形式化框架相比，其理论基础仍显不足。

\subsection{1.1. 组合性的维度}\label{dimensions-of-composability}

为了刻画动态组合的需求，我们在组合已被充分研究的代数性质之外，识别出两个正交的维度：

• 时间维组合性关注时间维度：当移除某个组件时，该组件对共享环境所做的修改必须被完全且安全地撤销。这要求跟踪组件执行过程中的每一次资源分配、事件注册和状态变更，并保证在移除时它们能被有序回收。

• 空间维组合性关注空间维度：组件必须能够以结构化、可验证的方式声明、发现并解析彼此之间的依赖关系。这要求管理依赖拓扑，并依据依赖变化协调各组件之间的生命周期。

在静态场景下，时间维组合性可以归结为词法作用域（例如 RAII {[}4{]}、括号模式 {[}5{]}），空间维组合性则可以归结为模块导入解析 {[}6{]}。而在动态场景下，组件在运行时不断加入和离开，两个维度都变得更加困难：时间维组合性必须处理作用域不受词法限定的、长期存续且带状态的效应；空间维组合性则必须处理在执行过程中出现、消失或身份发生变化的依赖关系。

\subsection{1.2. 动机示例}\label{motivating-examples}

\subsection{1.2.1. 插件系统}\label{plugin-systems}

插件系统是动态组合的典型实例。我们以 Visual Studio Code（VSCode）---最广泛使用的可扩展 IDE 之一---作为代表性例子。

时间维局限。VSCode 在名为扩展宿主（extension host）的共享进程中运行所有扩展。虽然扩展可以动态安装，但该宿主不提供在运行时卸载单个扩展代码的机制。一旦某个扩展的 activate 函数被执行，禁用或卸载它就需要重启整个宿主，从而影响到所有已加载的扩展。主题、快捷键绑定和代码片段等纯声明式扩展不携带代码，可以被自由移除。然而，在按安装量排序的前 100 个扩展中，有 87 个包含可执行代码1，因此在移除时需要进行这样的重启。尽管 VSCode 提供了 deactivate 钩子，但它只是在宿主进程终止时用作优雅关机的回调，因而并不能实现热移除（live removal）。此外，该钩子将效应清理与效应创建（位于 activate 中）相分离，违背了关注点局部性原则，使完整清理难以得到验证。

空间维局限。VSCode 确实提供了用于声明扩展间依赖关系的 extensionDependencies，但它很少被使用：在按安装量排序的前 100 个扩展中，只有 7 个对非内置扩展声明了 extensionDependencies。1 这种稀缺反映了扩展 API 的形态：它暴露的是命令、视图、语言特性等固定的、表层（surface-level）扩展点。扩展通过这些扩展点向宿主贡献能力，而不是彼此依赖，因此扩展间的依赖很少出现。此外，VSCode 的扩展间交互机制不提供任何结构性契约：它通过 vscode.extensions.getExtension(\ldots).exports 将某个扩展的功能暴露给其他扩展，但返回值是无类型的（默认为 any），因此依赖方无法依赖经过类型检查的接口。简而言之，VSCode 引导扩展朝向一组固定的、由宿主提供的扩展点，而不提供任何安全、结构化的方式让扩展彼此依赖。

这两个局限并非 VSCode 独有，而是普遍存在于各类插件系统之中 {[}2, 7{]}，只是程度有所不同。

\subsection{1.2.2. 自演化智能体框架}\label{self-evolving-agent-harnesses}

现代 AI 智能体依赖运行时的智能体框架 {[}8--10{]}。这类系统可以组合多样的工具套件 {[}11{]} 与执行环境，管理权限与沙箱，维护会话状态与持久化，提供上下文管理与记忆系统 {[}12{]}，编排子智能体与多智能体工作流 {[}13{]}，并向用户和自动化流程暴露接口。未来的智能体框架也许会在持续服务请求的同时，生成并部署对其自身组件的修改。模型合成的可复用工具为组件级的自我修改提供了范围更窄的先导 {[}14{]}。每一次这样的修改本身都是动态组合的一个实例。

由于这些修改持续发生，且只有有限甚至没有人工监督，动态组合性变得不可或缺。如果缺少时间维组合性，每一次自我修改都迫使系统完全重启，丢弃所有进程内累积的状态；在如此高的频率下，累计的不可用时间相当可观，进行中的任务也会被反复打断；更糟的是，一次有缺陷的自我修改甚至可能使本用于恢复的进程都无法运行。如果缺少空间维组合性，每个模块都必须自行检测并适配其所依赖模块的出现、消失或身份变化，而且只能通过临时手段（ad hoc）进行；更糟的是，朴素的代码替换策略可能悄然破坏依赖它的模块，或引入只有到重载时才会暴露的循环依赖。

\subsection{1.2.3. 粗粒度变通方案}\label{the-coarse-grained-workaround}

动态组合性之所以很少受到形式化研究的关注，一个原因是操作系统和容器编排器（container orchestrator）已经提供了粗粒度的替代方案。操作系统在进程粒度上提供时间维组合性；容器编排器 {[}3{]} 则在服务粒度上提供空间维组合性。在实践中，大多数软件通过求助于这些粗粒度机制来容忍细粒度组合性的缺失：行为异常的模块通过重启进程来处理，服务依赖则交给容器编排器管理。

然而，这种变通方案代价高昂。在时间维上，每次重启都会丢弃进程内累积的所有状态（例如缓存、连接、部分计算结果），而重建这些状态需要数秒到数分钟 {[}15{]}；在此期间维持可用性需要冗余副本，从而以额外资源开销来弥补无法恢复单个组件的不足。在空间维上，容器级编排无法表达共享地址空间的组件之间的依赖关系，并且会为本可作为本地函数调用的交互引入网络开销。这两种机制都作用在进程与容器的边界上，而现代系统却越来越倾向于在更细的粒度上进行组合。这种粒度不匹配要求一种与组件本身处于同一层级、用于管理效应和依赖的组合抽象。

1数据取自 Visual Studio Code Marketplace，检索时间为 2026 年 6 月 9 日。

\subsection{1.3. 主要贡献}\label{contributions}

动态组合性的两个维度分别关注计算如何修改其环境，以及计算如何依赖其环境。这正是效应系统 {[}16, 17{]} 与余效应系统 {[}18, 19{]} 所形式化的两个方向：效应为推理环境修改提供了形式化词汇，余效应则为推理环境需求提供了形式化词汇。然而，现有的形式化将推理限制在词法固定作用域上的编译期分析，并未扩展到组件在运行时加入与离开的动态场景。通过将效应提升为可逆的运行时模型、将余效应提升为响应式的依赖解析机制，我们得到了动态组合的统一形式化基础，它不依赖具体语言，可适用于任何需要动态组合的软件架构。我们做出如下贡献：

\begin{enumerate}
\def\labelenumi{\arabic{enumi}.}
\item
  我们将可逆效应形式化（第 3.1 节）：每一次上下文变换都带有一个由运行时跟踪的显式逆操作，跟踪与恢复都保持组合性，因此当组件被移除时上下文得以恢复。这建立了局部时间维组合性。
\item
  我们将响应式余效应形式化（第 3.2 节）：组件将其所需的余效应声明为规范，每当上下文发生变化时，系统即依据该规范，以激活、停用或中性三种状态之一通知组件。这建立了局部空间维组合性。
\item
  我们将效应上下文与余效应上下文统一为单一上下文类型（第 3.3 节），其中定义在余效应上的观测等价为效应提供了独立性，从而构成一种时空维组合性的编程范式。
\item
  我们给出动态组合演算（第 4 节）：它将这两种机制结合为组件的概念，并为组件的生命周期配备操作语义。其元理论将时空维组合性从单个组件推广到由交错组件构成的整个系统。
\item
  我们在 Cordis 中实现了上述思想（第 5 节）。Cordis 是一个时空维组合性的元框架，它提供一个实现形式化模型的核心库（包含效应跟踪与余效应解析），以及一个带配置对账和热模块替换的声明式组件加载器。
\end{enumerate}
